\section{Conclusions}
\label{sec:conclusions}

Reverse liquidity stress can be posed as the root of a continuous window function of a run-off multiplier, capped so that every rate stays a share, rounded up to six decimals, and refused with a reason when the target is not bracketed. On \kInst{} instances of a seeded corpus the engine implementing it agrees with an independent oracle on every outcome, returns a threshold on every solved instance (\kPThr{}\,\%) and answers none of the \kRefused{} refused ones, while a scan of the multiplier grid of the case design answers \kGFalse{}\,\% of them and overshoots the threshold by a median of \kGErrPc{}\,\%. Removing the cap, the refusal rule, the continuous window, the upward rounding or the effective capacity each breaks a different property.

On the bank of case C16 the multiplier that brings the survival period to 30 days is \mAdv{} on the Adverse base and \mExt{} on the Extreme base, and the 2023 fast-run base is refused because its ladder is shorter than the target. The fast-run base, calibrated only on rows dated before 2023, models \svbModel{} of the deposit base in the two days of the run against 0.85 observed, a ratio of \svbRatio{}, and the multiplier that reproduces the observation is \mSvb{} times the calibrated rate. The Adverse multiplier moves by \jointDevAdv{}\,\% over the joint band of conventions and is refused under the slow monetisation schedule. Nine of twelve predictions held; the hyperbolic rule, the stability over the joint band and the amber bar of the labelled run failed.

The next steps are to give the survival period a distribution through a stochastic depositor process and to report the multiplier as a quantile of that distribution, to extend the family beyond a common multiplier to segment-specific ones with a plausibility measure, and to repeat the labelled-run score on banks whose deposit behaviour is observed.
